% A4 quer mit zwei A5-Hälften: Text links, Notizen rechts.
\usepackage{fontspec}
\setmainfont{Latin Modern Roman}
\usepackage[ngerman]{babel}
\usepackage{amsmath,amssymb,amsthm,mathtools}
\usepackage{booktabs,tabularx,array}
\usepackage{tikz}
\usetikzlibrary{arrows.meta,calc}
\usepackage[a4paper,landscape,top=18mm,bottom=16mm,left=12mm,right=160.5mm]{geometry}
\usepackage{enumitem}
\makeatletter
\renewcommand\normalsize{\@setfontsize\normalsize{10.5}{13}}
\makeatother
\normalsize
\setlist{nosep,leftmargin=1.5em,topsep=.35em}
\pagestyle{empty}
\usepackage{microtype,needspace,etoolbox}
\usepackage[hidelinks]{hyperref}
\definecolor{cutsblue}{RGB}{22,89,112}
\definecolor{cutspale}{RGB}{231,241,244}
\newcommand{\N}{\mathbb N}
\newcommand{\Q}{\mathbb Q}
\newcommand{\R}{\mathbb R}
\newtheorem*{satz}{Satz}
\newtheorem*{definition}{Definition}
\theoremstyle{remark}
\newtheorem*{beispiel}{Beispiel}
\newcommand{\Beweisidee}{\par\smallskip\textbf{Beweisidee.}\ }
\setlength{\parindent}{0pt}
\setlength{\parskip}{.42em plus .08em minus .08em}
\setlength{\emergencystretch}{1.5em}
\setcounter{tocdepth}{1}
\newlength{\CutsWorkHeight}
\newcommand{\ExerciseSpace}[1]{%
 \par\smallskip\begingroup
 \setlength{\CutsWorkHeight}{\dimexpr\pagegoal-\pagetotal-2\baselineskip\relax}%
 \ifdim\CutsWorkHeight>#1mm\setlength{\CutsWorkHeight}{#1mm}\fi
 \ifdim\CutsWorkHeight>12mm
  \vfill{\small\sffamily Deine Rechnung oder Skizze}\par\nobreak
  \begin{tikzpicture}
   \draw[step=5mm,black!22,line width=.2pt] (0,0) grid (120mm,\CutsWorkHeight);
  \end{tikzpicture}
 \fi\endgroup}
\newcommand{\Exercise}[2]{%
 \clearpage
 \ifnum#1=1
  \section{Übungsblätter}\label{cuts:exercises}
  Arbeite zuerst ohne die Lösungen ab Seite~\pageref{cuts:solutions}.
  Nutze die Karoflächen für Skizzen und eigene Beweise.
 \fi
 \subsection*{Aufgabe #1: #2}\phantomsection\label{cuts:exercise-#1}}
\newcommand{\ExtraWork}[2]{%
 \clearpage\subsection*{Arbeitsraum zu Aufgabe #1}
 \textbf{Hinweis, falls du ihn brauchst.} #2\par\ExerciseSpace{130}}
\AddToHook{shipout/background}{%
 \begin{tikzpicture}[remember picture,overlay]
  \begin{scope}[shift={(current page.north west)},x=1mm,y=-1mm]
   \ifnum\value{page}>1
    \node[anchor=west,font=\sffamily\fontsize{7.5}{9}\selectfont,text=black!55]
       at (12,10) {BAND 19\quad DEDEKINDSCHE SCHNITTE\quad \CutsBookletLabel};
    \draw[black!25,line width=.25pt] (12,13)--(136.5,13);
    \node[font=\sffamily\fontsize{8}{10}\selectfont,text=black!60]
      at (74.25,201) {\CutsPageLabel{} \thepage};
   \fi
   \node[anchor=west,font=\sffamily\fontsize{8.5}{10}\selectfont,text=black!65]
      at (161,13) {NOTIZEN\quad Fragen, Ideen und Beweise};
   \begin{scope}[shift={(161,22)}]
    \draw[step=5mm,black!22,line width=.2pt] (0,0) grid (125,170);
   \end{scope}
   \node[font=\sffamily\fontsize{8}{10}\selectfont,text=black!60]
      at (223.5,201) {Notizen zu \CutsPageLabel{} \thepage};
  \end{scope}
 \end{tikzpicture}}
\newcommand{\CutsCoverFigure}{%
\begin{tikzpicture}[x=1cm,y=1cm]
 \fill[cutspale] (0,-.25) rectangle (6.2,.25);
 \draw[-{Stealth[length=2mm]},black!55] (0,0)--(10,0) node[right] {$\Q$};
 \draw[cutsblue,line width=1.3pt] (0,0)--(6.2,0);
 \draw[cutsblue,fill=white,line width=.9pt] (6.2,0) circle (2pt);
 \node[below=5mm] at (6.2,0) {$q$};
 \node[text=cutsblue] at (3.1,.8) {$\widehat q=\{p\in\Q\mid p<q\}$};
 \node[align=center,font=\small,text=black!65] at (5,-1.3)
 {Eine Zahl durch alles beschreiben,\\was links von ihr liegt.};
\end{tikzpicture}}
\allowdisplaybreaks[1]
\clubpenalty=10000
\widowpenalty=10000
\BeforeBeginEnvironment{satz}{\Needspace{6\baselineskip}}
\BeforeBeginEnvironment{definition}{\Needspace{6\baselineskip}}
\BeforeBeginEnvironment{beispiel}{\Needspace{5\baselineskip}}
\pretocmd{\section}{\Needspace{9\baselineskip}}{}{}
